\section{A dynamically changing sparse carrier}\label{sec:storage}

The previous theorems constrain motion and communication. We next obtain a
simultaneous storage condition with no assumption that the exceptional tape
has a fixed identity.

\begin{theorem}[Exact dynamic-carrier simulation]\label{thm:carrier}
\notformalized
Fix a native machine with $k\ge1$ tapes, $Q$ states and alphabet
$A\le2^a$. Fix integers $T,m\ge0$, put $n=2T+1$, and define
\begin{align*}
 V&=\sum_{j=1}^{\min(n,2m+2)}\binom nj(2^a-1)^j,
 &r&=\lceil\log_2(V+1)\rceil,\\
 B_0&=k(T+1)n^{k+1}(n+1)^k(4T+1)(1+n^2).
\end{align*}
There is one native machine with one work tape and the same alphabet that,
on every source input halting within $T$ for which every reached configuration
has at most one tape with more than $m$ nonblank cells, reproduces exactly the
source's successful-transition count, raw-head moves and append times/bits.
It has at most $1+QB_0 2^{(k+1)r}$ states. It may fail or halt off the promise.
The table depends on the source machine, $T,m$, not on the individual input.
\end{theorem}
\begin{informalproof}
Embed the alphabet in $\F_2^a$, with blank zero. A tape supported in
$[-T,T]$ is a vector with $n$ blocks. There are at most $V$ nonzero vectors of block
weight at most $2m+2$. A uniformly chosen binary $r\times an$ matrix kills
each with probability $2^{-r}$, so some matrix $H$ kills none. Fix it.
Its syndrome is injective on legal tape words of weight at most $m+1$.
The partial inverse is a finite transition-table function; it does not
assert that a dense tape can be decoded. This finite selection is checked in
\lean{binary_matrix_collision_card} and \lean{exists_sparse_injective_hash};
\lean{binary_hash_bits_eq_ceiling} identifies the bit-length convention.
The remaining native transition table and handoff invariant below are not
yet formalized.

Maintain all $k$ full-tape syndromes and exact nonblank counts in control.
Replacing old symbol $u$ by $v$ at address $z$ updates the syndrome by
$H(e_z\otimes(\operatorname{enc}(u)+\operatorname{enc}(v)))$ and the count
by $1_{v\ne0}-1_{u\ne0}$. The syndrome identity is the specialization of
\lean{LeanMachine.linearHash_update} to the alphabet embedding.
Store all logical positions, the source state and
time, an active tape $f$, physical position $p$, offset $u_0$ with
$p+u_0=y_f$, a frozen phase-entry syndrome, and an initially empty interval
$I$ of physical cells previously scanned in this phase.

The invariant before a source step is: every inactive tape is $m$-sparse;
the physical symbol at $j\in I$ is the current logical symbol of $f$ at
$j+u_0$; outside $I$ the logical symbol is given by the frozen phase-entry
word, whose weight is at most $m+1$. The syndrome registers always encode
the current whole logical tapes, including a dense active tape. Initially
all words and registers are zero, $f=0$, $u_0=0$ and $I=\varnothing$.

Read the old inactive symbols by sparse decoding. Read the active symbol
physically if $p\in I$; otherwise ignore the stale physical symbol and use
the frozen snapshot. This gives the entire original old tuple. Apply the
source transition once. If it halts, halt without a successful step. Otherwise
update every logical syndrome, count and position. Physically write the old
active tape's prescribed new symbol at $p$ and take its prescribed unit move.
Take the original raw move and optional append in this same transition.

If the updated configuration has no dense tape, or its dense tape is still
$f$, retain the carrier, offset and snapshot, and replace $I$ by $I\cup\{p\}$.
Unit motion makes this a single interval; its new cell was just written to
the correct value. If a different tape $g$ has become the unique dense tape,
it was inactive and $m$-sparse before the transition. One write raises its
count by at most one, so its new count is exactly $m+1$. Decode its updated
syndrome as the new frozen snapshot, choose the new offset
$u_0'=y_g'-p'$, and reset $I$ to empty. The old active tape is now sparse
and can be decoded from its updated live syndrome. The physical move was
still the old carrier's move; the new offset requires no second move, seek,
clearing scan or copy. Stale physical cells are ignored until first visited
and overwritten. This proves the invariant through arbitrary handoffs,
including simultaneous changes on both involved tapes.

All choices are finite table operations. Illegal decodes, multiple dense
updated counters and out-of-range fields may halt; promised executions never
take these branches. Force a halt when the stored time is $T$. A promised
source already halts there or earlier. The physical and logical positions
are bounded by $T$, the offset by $2T$, and $I$ has at most $1+n^2$ possible
empty/endpoint descriptions. Time, carrier, positions, counts and offset
have exactly the factors overcounted in $B_0$. There are $k$ live syndromes
and one frozen syndrome. Multiplying by $Q$ and adding a default state gives
the stated bound. This is an actual finite native table with one physical
work head; it is nonuniform in $T,m$, and no small table-description claim
is needed.
\end{informalproof}

\begin{theorem}[Theorem S: simultaneous large stores]\label{thm:stores}
\partformal{StoresFromCarrier.lean}{stores_tail_of_carrier,native_two_dense_heads,native_support_window_complete}{The exact state-bit budget and success tail are checked for one common native carrier simulator with the stated state count and halt/output correspondence. Constructing that simulator, as in \cref{thm:carrier}, remains unformalized.}
Fix $k,a,C$. There are $\delta,\eta>0$ and $d_0$, uniform over all $b\ge1$,
such that fewer than $2^{(1-\eta)P}$ timely correct inputs fail to reach a
configuration with two distinct tapes each containing more than $\delta P$
nonblank cells. This holds for fixed controllers and, more generally, for
$q\le P/(128N_1^2)$, where $N_1=\lfloor160aC^2\rfloor+1$.
\end{theorem}
\begin{informalproof}
Apply \cref{thm:carrier} to precisely the successful inputs satisfying the
one-dense-tape promise, with $T=\lfloor C(P+1)\rfloor$ and
$m=\lfloor\delta P\rfloor$. The one-head case of \cref{thm:returns} applies
to each nonuniform simulator if its state-bit cost is at most
$\rho P$, $\rho=1/(64N_1^2)$. Write $K=4C+1$. At large $P$,
$2m+2\le4\delta P\le n/2$ and
\[
 r\le1+4\delta P\log_2\frac{eK(2^a-1)}{4\delta}.
\]
Choose $0<\delta\le1/8$ so that the last coefficient multiplied by $k+1$
is at most $\rho/4$. The logarithm of $B_0$ is $O(\log(P+1))$. The source
state cost is at most $\rho P/2$, leaving sufficient room for this overhead
and ceilings. The checked \lean{carrier_state_bits_small} proves the complete
integer bit budget directly from the sparse-menu generating-function bound,
including the default state and every syndrome. Thus the simulator satisfies
the one-head budget and has fewer
than $2^{(1-1/(8N_1))P}$ timely correct inputs. Every promised successful
source input is among them, with no loss and no per-input simulator choice.
All thresholds depend only on the stated resources, because $P\ge d^2$.
The checked complement lemma converts the floor cutoff to the strict real
threshold at one common reached configuration, including the halt configuration.
The finite window $[-T,T]$ contains the entire native nonblank support.
\end{informalproof}

\begin{lemma}[Persistent unchanged cells]\label{lem:persistent-cells}
\formal{Storage.lean}{native_support_growth_le_one,native_support_shrink_le_one,ExecutionPrefix.untouched_cell_constant,ExecutionPrefix.persistent_cells}
For any fixed finite window on a tape, one native transition changes its
nonblank-cell count by at most one. In an interval of $\ell$ successful steps,
at most $\ell$ terminal nonblank cell addresses in that window were scanned
in the interval; every other terminal nonblank cell in the window keeps its
symbol throughout the interval.
\end{lemma}
\begin{outline}
Outside the old scanned address, \cref{lem:old-write} preserves the symbol.
The touched set is the image of the interval's times under the head-position
map and has cardinality at most its duration. Remove it from the terminal
support. The formal statement uses any finite window containing the cells;
all live cells are in $[-T,T]$ on a $T$-step run.
\end{outline}

At the witness time of \cref{thm:stores}, both counts exceed $\delta P$.
They started at zero, so the time exceeds $\delta P$. Throughout the preceding
$\lfloor\delta P/2\rfloor$ steps both counts exceed $\delta P/2$.
Deleting the addresses written in that interval leaves more than
$\delta P/2$ fixed nonblank cells on each tape. Their values need not be
independent, unknown, necessary for output, or even read again.
The checked \lean{ExecutionPrefix.dense_persistent_interval} gives the more
general exact loss of $\ell$ cells over a preceding interval of $\ell$ steps.
